\documentclass[runningheads]{llncs}
\usepackage{graphicx}
\usepackage{amsmath,amssymb,amsfonts}
\usepackage{booktabs}
\usepackage{array}
% \usepackage{microtype}
\usepackage{algorithm}
\usepackage{algpseudocode}
\usepackage{xurl}
\setlength{\textfloatsep}{6pt plus 1pt minus 1pt}
\setlength{\floatsep}{5pt plus 1pt minus 1pt}
\setlength{\intextsep}{5pt plus 1pt minus 1pt}
\sloppy
\emergencystretch=3em

\begin{document}

\title{Policy-Aware Lexicographic CP-SAT for Sequential Constrained Subset Selection in Examination Assembly}
\titlerunning{Policy-Aware Lexicographic CP-SAT}

\author{Anonymous Author(s)}
\authorrunning{Anon. Author(s)}
\institute{Institute Anonymized for Double-Blind Review\\
\email{author@domain.org}}

\maketitle

\begin{abstract}
Sequential examination-paper assembly is a repeated constrained subset-selection problem with exact resources, pairwise exclusions, distributional targets, and sequential carry-over decisions. We audit whether sequential lexicographic decomposition---solving sections independently while carrying selected items forward---preserves joint whole-paper quality. The study compares sequential lexicographic, joint lexicographic, and five alternative CP-SAT policies on a production-derived 475-item CBSE Class X Mathematics corpus across five independently generated 80-blueprint manifests. On the fixed manifest, sequential $B_1$ and joint $B_J$ tie on all 70 common-feasible cases; four section orderings also retain identical aggregate medians. Across manifests, $B_1$ has seed-level median difficulty deviation 38--40 versus 50--52 for first-feasible $B_0$. Direct-compliance and weighted scalar controls match $B_1$'s median difficulty with lower Bloom deviation and faster runtime, making them practical when explicit priority governance is unnecessary. Bounded lookahead ties $B_1$ on quality while raising median runtime from 0.124 to 0.257 seconds. Results characterize when sequential decomposition is safe in the tested regime and provide a reproducible benchmark for policy-layer optimization in constrained subset-selection workloads.
\keywords{Constraint Programming \and CP-SAT \and Lexicographic Optimization \and Multi-Objective Optimization \and Constrained Subset Selection \and Automated Test Assembly.}
\end{abstract}

%===============================================================================
\section{Introduction}
\label{sec:intro}
%===============================================================================

Large-scale assessment platforms manage item repositories subject to exact marks and duration, section counts, Bloom-level balance~\cite{bloom1956,anderson2001}, chapter coverage, and duplicate avoidance. India's National Education Policy~\cite{moenep2020} further emphasizes competency-based assessment. More generally, this is repeated constrained subset selection with exact resource constraints, pairwise exclusions, distributional targets, and sequential carry-over decisions.

\subsection{The Decomposition Question}
Production systems often solve papers section by section, carrying selected items forward to enforce uniqueness. This decomposition is modular, auditable, and scalable, but risks sequential lock-in: an early optimum may consume items needed later for global quality. For an optimization workflow, this is a practical decomposition question rather than only an implementation detail: the modeller must know whether smaller sequential subproblems preserve the outcome of a larger joint model.

The central algorithmic question is therefore whether sequential lexicographic decomposition preserves the quality of joint whole-paper lexicographic optimization, and when a deployment engineer can adopt sequential decomposition without measurable quality loss.

This question is distinct from objective-function choice in ATA. Prior work~\cite{vanderLinden2005,theunissen1985,ignjatovic2021,li2021,nguyen2021} mainly optimizes a single assembled form or studies heuristic/metaheuristic search. We instead compare deployment-oriented policy layers over a shared CP-SAT model, with sequential--joint equivalence as the primary audit.

A secondary structural question is the feasibility gap: a feasible paper need not align with blueprint quality targets. A first-feasible solver ($B_0$) may return any valid assignment, while optimized $B_1$ aligns the difficulty histogram as a primary objective. The size of this gap, and which policy closes it most efficiently, motivates the research questions below.
\subsection{Research Questions and Scope}
This work investigates four questions in a sequential generation setting:
\begin{itemize}\setlength{\itemsep}{1pt}\setlength{\parskip}{0pt}
\item \textbf{RQ1 (Structural Alignment):} How much does multi-objective optimization improve metadata-histogram alignment over first-feasible generation ($B_0$), with chapter coverage reported as a saturation check?
\item \textbf{RQ2 (Solver Runtime):} What solver-runtime overhead does optimization introduce as item pools scale to $N=10{,}000$, excluding pre-solver graph construction?
\item \textbf{RQ3 (Blueprint Tightness):} How does optimization gain behave as the feasible region contracts under tighter chapter-eligibility restrictions?
\item \textbf{RQ4 (Soft Compliance):} To what extent does explicit multi-pass optimization improve soft blueprint compliance metrics over first-feasible generation?
\end{itemize}
\paragraph{Scope and boundaries.} Duplicate detection constructs a conflict graph $S$ supplying at-most-one constraints; it is a pipeline boundary condition, not an in-solver objective. Evaluation measures structural alignment and solver-only runtime; embedding and graph-construction costs are excluded, and psychometric IRT validity is not evaluated.
\paragraph{Contributions.} The paper makes three algorithm-engineering contributions: (i) a decomposition audit showing when sequential lexicographic CP-SAT matches joint whole-paper optimization; (ii) a comparison of first-feasible, scalar, weighted, lexicographic, joint, and lookahead policies under a common CP-SAT model; and (iii) a reproducible constrained subset-selection benchmark with manifests, fixtures, and result artifacts.
\section{Related Work}
\label{sec:related}
Automated Test Assembly is commonly formulated as 0--1 integer programming under content and test-information constraints~\cite{theunissen1985,vanderLinden2005,vanderLinden1989,lord1980}, with later work addressing parallel forms, cognitive-diagnosis settings, and multiobjective assembly~\cite{ignjatovic2021,li2021,nguyen2021}. Unlike IRT-centred ATA, this study uses operational metadata---marks, time, difficulty, chapter, and Bloom labels---for repositories without full psychometric calibration. Weighted scalarization and lexicographic priorities are established multiobjective approaches~\cite{ehrgott2005,ignizio1976}, and CP-SAT supports the Boolean, linear, and conflict constraints required here~\cite{perron2023,nieuwenhuis2006}. The contribution is a policy-layer comparison of sequential lexicographic decomposition, joint optimization, scalar controls, and bounded lookahead under one reproducible CP-SAT benchmark, framing ATA as repeated constrained subset selection.
\section{Problem Formalization}
\label{sec:formalization}
Let $Q=\{q_1,\dots,q_n\}$ be an item bank. Each item $q_i=(m_i,t_i,T_i,d_i,b_i)$ has marks $m_i\in\mathbb{Z}^+$, completion time $t_i\in\mathbb{Z}^+$, topic tags $T_i\subseteq\mathcal{T}$, difficulty band $d_i\in\{1,\dots,|D|\}$, and Bloom's cognitive level $b_i\in\mathcal{B}$. An eligibility set $E\subseteq\mathcal{T}$ filters candidate items: only items satisfying $T_i\cap E\neq\emptyset$ are included in the solver pool. A required-coverage set $R\subseteq\mathcal{T}$ specifies chapters to be represented. In this paper, \emph{chapter coverage} is the instantiated dataset metric $\sum_{t\in R}y_t$, i.e., the number of required chapters represented by at least one selected item; the term is used consistently for this quantity.
A blueprint is specified as $B=(M,W,K,E,R,\Delta,\mathcal{H},\mathcal{S})$, where $M$ is total marks, $W$ is duration, $K$ is total question count, $\Delta=(\Delta_1,\dots,\Delta_{|D|})$ is the target difficulty distribution ($\sum_k\Delta_k=K$), and $\mathcal{H}$ and $\mathcal{S}$ represent hard and soft blueprint constraints.
\paragraph{Decision variables and baseline $B_0$.}
Let binary decision variable $x_i\in\{0,1\}$ represent item selection ($x_i=1$ iff $q_i$ is selected). The baseline first-feasible solver ($B_0$) finds $x\in\{0,1\}^n$ satisfying:
\begin{align}
\sum_{i=1}^n m_i x_i &= M \quad &&\text{(Exact total marks)} \label{eq:marks} \\
\sum_{i=1}^n t_i x_i &= W \quad &&\text{(Exact total duration)} \label{eq:time} \\
\sum_{i=1}^n x_i &= K \quad &&\text{(Exact question count)} \label{eq:count} \\
x_i + x_j &\le 1 && \forall (i,j)\in S \quad\text{(Anti-duplicate conflicts)} \label{eq:conflict} \\
x_i &= 0 && \forall i:\,T_i\cap E=\emptyset \quad\text{(Topic eligibility)} \label{eq:eligibility}
\end{align}
where $S\subseteq Q\times Q$ is the pre-computed set of conflicting near-duplicate pairs. Required chapter coverage is modeled using binary indicators $y_t\in\{0,1\}$ for $t\in R$:
\begin{equation}
y_t \le \sum_{i:\,t\in T_i} x_i \quad \forall t\in R \label{eq:coverage_link}
\end{equation}
For policies that optimize coverage, this upper-bound linking is sufficient because maximizing $y_t$ forces $y_t=1$ whenever chapter $t$ is represented. In the reported manifest, coverage is an objective indicator rather than a separate hard requirement; it is not imposed as a mandatory lower bound. This distinction matters because coverage can be optimized, reported, or saturated without changing hard feasibility. For policies that do not optimize coverage, reported coverage is computed directly from the selected item topics rather than from unconstrained $y_t$ values. The baseline $B_0$ has no objective function; it halts at the first assignment satisfying the hard constraints.
This modelling choice also reflects the production setting: infeasibility should be driven by hard resource, eligibility, and conflict constraints, while representational breadth is treated as a preference to be optimized when available capacity permits.
\paragraph{Algorithmic properties.}
For one section with $n_s$ eligible items, $e_s$ conflict edges, and $h_s$ additional hard rules, the base model has $n_s$ item variables, three aggregate equalities, $e_s$ conflict inequalities, and the additional rules. In the seven-band, three-Bloom implementation, $B_1$ adds deviation variables, coverage indicators, and Bloom-deviation variables, with standard absolute-value and coverage-link constraints. These counts explain policy overhead rather than give a universal model-size formula. The three-pass policy invokes CP-SAT up to three times per section under a 30\,s cap.
The feasibility problem is NP-hard because cardinality-constrained subset sum is a special case: assign item marks to the input integers, fix count to the required cardinality, set all times to one with $W=K$, and omit conflicts and optional rules. The exact-marks constraint then asks whether a size-$K$ subset sums to $M$~\cite{garey1979}. The count--time dynamic programme used by $B_{LS}$ is therefore pseudo-polynomial in the integer time bound, not a polynomial-time solution to the general feasibility problem.
\paragraph{Sequential decomposition.}
Let $X_s$ be the items selected before section $s$. Sequential lexicographic optimization is equivalent to whole-paper lexicographic optimization only when each section optimum leaves a jointly optimal completion for later sections and no coupling exists except through $X_s$. This condition is strong: it can fail through early item consumption, scarce high-value items, or cross-section soft-objective interaction. Thus $B_1$ is not claimed lossless in general; $B_J$ tests this empirically and provides the main decomposition control.
%===============================================================================
\section{Policy Formulations and Optimization Models}
\label{sec:proposed}
%===============================================================================
We evaluate policy implementations rather than a new solver algorithm. The core comparison contrasts first-feasible search with an independent constructive--local baseline, two single-objective controls, and sequential and joint lexicographic CP-SAT policies under the same blueprint constraints. This allows differences in outcomes to be attributed to policy design rather than to changes in the underlying hard feasibility model.
\subsection{Bounded Lexicographic CP-SAT Policy ($B_1$)}
$B_1$ implements a three-pass lexicographic CP-SAT policy. When every pass is \texttt{OPTIMAL}, it realizes the stated priority order; a time-capped \texttt{FEASIBLE} pass carries only an incumbent bound:
\begin{enumerate}\setlength{\itemsep}{1pt}\setlength{\parskip}{0pt}
    \item \textbf{Pass 1 (Primary Difficulty Alignment):} Minimize $L_1$ difficulty deviation $D(x) = \sum_{k=1}^{|D|} u_k$, where auxiliary integer variables $u_k \ge 0$ linearize the absolute error:
    \begin{align}
    u_k &\ge \sum_{i:\,d_i=k} x_i - \Delta_k && \forall k, \label{eq:u_bounds} \\
    u_k &\ge \Delta_k - \sum_{i:\,d_i=k} x_i && \forall k.
    \end{align}
    When Pass~1 is \texttt{OPTIMAL}, let $D^*$ be its optimal deviation; for a time-capped \texttt{FEASIBLE} result, $D^*$ denotes the incumbent deviation.
    \item \textbf{Pass 2 (Secondary Chapter Coverage):} Fix the Pass~1 difficulty value by enforcing $\sum_k u_k \le D^*$, then maximize required-chapter coverage $\sum_{t\in R} y_t$. Let $C^*$ denote the optimal coverage score, or the incumbent score if the pass ends with a time-capped \texttt{FEASIBLE} result.
    \item \textbf{Pass 3 (Tertiary Bloom Fit):} Fix the previous objective values with $\sum_k u_k \le D^*$ and $\sum_{t\in R}y_t=C^*$, then minimize Bloom deviation $P(x)=\sum_{g\in\mathcal{B}} |\sum_{i:b_i\in g}x_i-T_g|$. Here $T_g$ is the target count for cognitive group $g$.
\end{enumerate}
Each pass uses a fresh CP-SAT model with prior objective values fixed as constraints; selected item IDs are carried to the next section. The design deliberately exposes the priority order instead of relying on a single aggregate score, making it easier to explain why one feasible paper is preferred over another.
\noindent\textbf{Status handling.} The 30\,s budget is a safeguard; all feasible RQ1/RQ4 runs completed below 2.0\,s. A \texttt{FEASIBLE} timeout satisfies hard constraints but does not certify optimality. All reported RQ1 common-feasible $B_1$ passes were \texttt{OPTIMAL}: 1,050/1,050 pass-level statuses across 70 papers (artifact: \texttt{results\_b1\_status\_audit.json}).
\subsection{Alternative Baseline and Control Policies}
\begin{itemize}\setlength{\itemsep}{1pt}\setlength{\parskip}{0pt}
    \item \textbf{$B_{LS}$ (Constructive--Local Baseline):} A non-CP-SAT constructive dynamic programme selects an exact count--time seed using greedy metadata priorities, followed by deterministic improving equal-time swaps. The dynamic programme enforces hard feasibility; its scoring and local search do not optimize the CP-SAT objectives. It provides a useful independent baseline because it avoids CP-SAT modelling choices while still attempting metadata-aware improvement.
    \item \textbf{$B_2$ (Weighted-Sum Scalarization):} Solves a single CP-SAT pass minimizing $\alpha D(x) + \beta U(x) + \gamma P(x)$, where $U(x) = |R| - \sum_{t\in R} y_t$. The default shared-suite weight vector is $(\alpha,\beta,\gamma)=(1,1,1)$, using raw rather than normalized objective components; results should therefore be interpreted as specific to this policy scale.
    \item \textbf{$B_C$ (Direct Compliance Control):} A joint whole-paper CP-SAT model minimizes $D_{\mathrm{total}}+P_{\mathrm{paper}}$, the affine equivalent of the reported complete-paper soft-compliance loss because every template paper has 38 questions. It is a single-objective control, not a new optimization method.
    \item \textbf{$B_J$ (Joint Lexicographic Control):} A whole-paper CP-SAT model applies the same $D\!\to\!C\!\to\!P$ priority order as $B_1$, while enforcing all section constraints and cross-section uniqueness simultaneously.
\end{itemize}
\subsection{Bounded Lookahead-Guided Sequential Diagnostic}
To test whether sequential lock-in can be reduced without changing the decomposition, bounded lookahead generates up to $k=3$ distinct $B_1$ candidates per section and scores downstream states using a coarse relaxed estimator. Let $P_s(x)$ be the remaining-section pool after excluding candidate $x$, and let $r_s,w_s$ be its required count and time:
\begin{align}
R_s(x) &= \mathbf{1}\left\{\sum_{q\in P_s(x)}1\ge r_s\right\}, \\
\sigma_s(x) &= \left|\sum_{q\in P_s(x)}t_q-w_s\right|, \\
S(x) &= 100\sum_s R_s(x)-\sum_s \sigma_s(x).
\end{align}
The highest-scoring candidate is selected. The estimator is diagnostic, not a completion counter, and does not certify future feasibility. Each candidate-generation subsolve uses one worker, seed 1, and a 5\,s cap. The lookahead control therefore tests whether a small amount of downstream information can reduce sequential lock-in without replacing the production decomposition.
\subsection{Constructive--Local Baseline Details}
The constructive--local baseline uses a count--time dynamic program to build a feasible section, then applies deterministic equal-time swaps that improve $(D,P,-C)$ while preserving conflicts. Its state complexity is $O(n_sKW)$ without conflicts, with an additional implementation factor for conflict checks. Unlike CP-SAT policies, it is heuristic and provides no optimality certificate.
\subsection{Compliance Metrics}
Blueprint compliance is partitioned into non-negotiable hard rules $\mathcal{H}$ (e.g., exact marks, time, question count) and soft distributional goals $\mathcal{S}$ (difficulty alignment and Bloom fit). For hard compliance, $CR_{\mathrm{hard}}(x) = 1.0$ for any valid output. Soft compliance is scored as:
\begin{equation}
\begin{aligned}
CR_{\mathrm{soft}}(x) &= \frac{1}{2}\left(1-\min\left\{1,\frac{D(x)}{2K}\right\}\right) \\
&\quad + \frac{1}{2}\left(1-\min\left\{1,\frac{P(x)}{2K}\right\}\right).
\end{aligned}
\end{equation}
For a section, $D=\sum_{k=1}^{7}|\sum_{i:d_i=k}x_i-\Delta_k|$ denotes difficulty deviation and $P=\sum_{g\in\{\mathrm{RU,APP,AEC}\}}|\sum_{i:b_i\in g}x_i-T_g|$ denotes Bloom-group deviation. Because observed and target histograms both sum to $K$, both are bounded by $2K$.
For a complete paper, $K_{\mathrm{tot}}=38$, $D_{\mathrm{total}}$ sums section-level difficulty deviations, and $P_{\mathrm{paper}}$ uses the 54/24/22 Bloom target. Thus $0\leq D_{\mathrm{total}},P_{\mathrm{paper}}\leq76$, clipping is inactive, and $CR_{\mathrm{soft}}=1-(D_{\mathrm{total}}+P_{\mathrm{paper}})/(4K_{\mathrm{tot}})$.
The weighted policy $B_2$ uses $U=|R|-\sum_{t\in R}y_t\in[0,7]$ and the fixed raw objective $D+U+P$. Its unnormalized weights were specified before the experiments rather than tuned from the results. Hence $B_2$ represents a scalar preference model, not the same priority semantics as lexicographic $B_1$.
%===============================================================================
\section{Experimental Evaluation}
\label{sec:experiments}
%===============================================================================

\subsection{Dataset and Benchmark Setup}
The empirical corpus contains \textbf{475 CBSE Class X Mathematics items} from a production assessment system, each with marks, completion time, difficulty band, Bloom tag, and chapter metadata. The bank is easy-skewed (83\% in bands 1--3) and marks-skewed (73\% 1-mark items), supporting a production-derived test of difficulty alignment under realistic imbalance. For RQ2, synthetic pools calibrated to empirical marginals are generated up to $N{=}10{,}000$; conflict density and feature correlations are not claimed to reproduce the production bank. The synthetic experiment is therefore a solver-scaling diagnostic, not evidence of population-wide corpus realism.
\subsection{Blueprint Manifest Generation}
For RQ1, 80 blueprints are generated deterministically from \texttt{random.Random(42)}. The corpus defines a fixed seven-chapter universe; each blueprint samples 5--7 eligible chapters, draws difficulty targets using weights $(4,5,4,3,2,1,0.3)$ with cyclic remainder assignment, and uses the 54/24/22 paper-level Bloom split.
Each blueprint has five sections (A--E), together containing 38 questions. A/B/C/D/E contain 20/5/6/4/3 questions, 20/10/18/20/12 marks, and 40/15/36/40/18 minutes, respectively. Section-level targets are stored in the manifest; sections are solved sequentially A$\to$E while carrying selected IDs forward.
\subsection{Whole-Paper Control on Excluded Blueprints}
A whole-paper joint CP-SAT audit classifies all 10 sequentially excluded blueprints as \texttt{INFEASIBLE} under the same hard constraints, indicating global blueprint infeasibility rather than sequential lock-in; the common-feasible baseline is therefore 70/80. This audit is important because it prevents infeasible manifest rows from being misinterpreted as failures of the sequential policy.
\subsection{Results}
\paragraph{RQ1 (Structural Alignment).}
The fixed manifest provides paired policy comparisons. On the fixed manifest, $B_1$ improves $D$ over $B_0$ in all 70 pairs (median difference 12, 95\% CI $[10,12]$, Holm-adjusted $p=3.6{\times}10^{-12}$, $r=0.87$). These are paired results for the tested manifest, not independent evidence from 70 unrelated experiments. Against $B_{LS}$, $B_1$ improves $D$ in 22 pairs and ties in 48 (median difference 0, CI $[0,0]$, adjusted $p=1.4{\times}10^{-5}$, $r=0.56$). This establishes structural-alignment differences only, not psychometric or pedagogical superiority.
\begin{table}[t]
\caption{RQ1 outcomes on the 70 blueprints feasible for all six core policies. Entries are medians with IQRs, except coverage, which is a median count. $D$ is summed section-level difficulty deviation; Bloom is complete-paper $L_1$ deviation from the 54/24/22 split. Runtime is complete-paper model construction plus solving; the 10 exclusions are jointly infeasible under hard constraints.}
\label{tab:rq1}
\centering
\small
\begin{tabular}{@{}lccccc@{}}
\toprule
\textbf{Model} & \textbf{Feasible} & $D\downarrow$ & \textbf{Coverage} & \textbf{Bloom}$\downarrow$ & \textbf{Solver s} \\
\midrule
$B_0$ first-feasible        & 70/80 & 52 [48,56] & 6 & 40 [40,40] & 0.03 \\
$B_{LS}$ constructive--local & 70/80 & 42 [38,44] & 6 & \textbf{22} [22,26] & 0.73 \\
$B_2$ weighted-sum          & 70/80 & \textbf{40} [36.5,44] & 6 & \textbf{22} [22,26] & 0.07 \\
$B_C$ direct compliance    & 70/80 & \textbf{40} [38,44] & 6 & \textbf{22} [22,24] & \textbf{0.04} \\
$B_J$ joint lexicographic  & 70/80 & \textbf{40} [36,44] & 6 & 24 [22,26] & 0.17 \\
$B_1$ sequential lexicographic & 70/80 & \textbf{40} [36,44] & 6 & 24 [22,26] & 0.13 \\
\bottomrule
\end{tabular}
\end{table}
$B_C$ and $B_2$ match $B_1$'s median difficulty while giving lower Bloom deviation and runtime. Because clipping is inactive, $CR_{\mathrm{soft}}$ is an affine transform of $D_{\mathrm{total}}+P_{\mathrm{paper}}$, making $B_C$ affine-equivalent to the reported loss but not to lexicographic $B_1$. Thus $B_C$ and $B_2$ are attractive scalar controls, while $B_1$ preserves an explicit, weight-free priority order. The comparison therefore separates two useful deployment choices: optimize a scalar compliance score when latency dominates, or use lexicographic control when stakeholders require inspectable priority semantics. Coverage is weakly discriminating here: paired core-policy comparisons tie in represented-chapter count.
\paragraph{Baseline sensitivity.} $B_0$ uses one worker, deterministic search, and retained per-blueprint seeds. $B_2$ uses the pre-specified raw $(1,1,1)$ weights. A grid over $\{1,3,10\}^3$ retains median coverage 6 and yields median $D\in\{40,41\}$, Bloom deviation $\in\{22,23,24\}$, and runtime $0.072$--$0.126$\,s on the same 70 rows (artifact: \texttt{results\_weight\_sensitivity.json}). This shows policy-scale sensitivity, not tuned weights.
The bounded-lookahead diagnostic ties sequential $B_1$ on all quality metrics across 70 comparable papers, while increasing median runtime from $0.124$ to $0.257$\,s; it therefore gives no quality gain on this corpus.
\begin{table}[t]
\caption{Bounded-lookahead diagnostic on the fixed manifest. Entries are medians with IQRs. Positive quality differences favour lookahead; runtime difference is Sequential $B_1$ minus Lookahead.}
\label{tab:lookahead}
\centering
\small
\begin{tabular}{@{}lccc@{}}
\toprule
Metric & Sequential $B_1$ & Lookahead $k{=}3$ & Difference ($B_1-$lookahead) \\
\midrule
Difficulty $D$ & 40 [36,44] & 40 [36,44] & 0 \\
Bloom deviation & 24 [22,26] & 24 [22,26] & 0 \\
Chapter coverage & 6 [6,7] & 6 [6,7] & 0 \\
Solver runtime (s) & 0.124 [0.109,0.140] & 0.257 [0.226,0.288] & -0.132 \\
\bottomrule
\end{tabular}
\end{table}
\paragraph{Robustness, Joint Control, and Ablations.}
Robustness is assessed across five manifests (seeds 42--46), treating the manifest as the resampling unit. Across clusters, $B_1$ has seed-level median $D$ values of 38--40 versus 50--52 for $B_0$. Paired effects, bootstrap intervals, ties, and Holm-adjusted tests are in \texttt{results\_multiseed\_revision.json}; pooled rows are descriptive only. This cluster-level treatment avoids overstating evidence from many blueprint rows generated within the same manifest.
On the fixed manifest, $B_J$ and sequential $B_1$ tie on difficulty, Bloom deviation, and coverage in all 70 jointly feasible cases; $B_J$ is slower. Alternative section orders also give identical aggregate medians. Objective ablations yield median $(D,\mathrm{Bloom})$ of $(40,32)$ for D-only, $(40,38)$ for D$\to$C, $(40,24)$ for D$\to$C$\to$P, and $(48,22)$ for C$\to$P; chapter coverage again has median 6.
The joint-control result is especially relevant for optimization design because it separates decomposition risk from objective semantics. In the tested regime, the sequential policy behaves like its joint counterpart despite solving smaller section-level models; however, the ablation results show that the objective order still matters, particularly for Bloom fit.
\paragraph{RQ2 (Solver Scaling).}
For Section A ($K{=}20$), median runtime rises from $0.003$/$0.021$ s for $B_0$/$B_1$ at $N=100$ to $0.326$/$1.643$ s at $N=10{,}000$. Thus $B_1$ incurs a $2$--$7\times$ overhead while remaining within budget. A descriptive log--log fit gives slopes of about $1.00$ for $B_0$ and $0.93$ for $B_1$; these are empirical, not asymptotic, slopes.
The scaling result should therefore be read as evidence that the additional lexicographic passes are affordable for the tested section size and synthetic pool range, not as a claim that larger templates or denser conflict graphs will preserve the same slope.
\paragraph{RQ3 (Blueprint Tightness).}
In a 100-case tightness sweep, reducing eligible chapters from 7 to 3 lowers common $B_0$/$B_1$ feasibility from 20/20 to 5/20. Conditional on common feasibility, $B_1$ maintains a median difficulty improvement of 6--12 units over $B_0$. The legacy runner records common-feasible versus non-common-feasible cases, not policy-specific \texttt{INFEASIBLE}/\texttt{UNKNOWN} categories.
\paragraph{RQ4 (Soft Compliance).}
Across 66 common-feasible pairs, complete-paper $CR_{\mathrm{soft}}$ rises from median $0.408$ to $0.566$, and $B_1$ outperforms $B_0$ in all cases (Wilcoxon $p=1.6{\times}10^{-12}$, $r=0.87$). This should be read with tier-wise feasibility counts rather than as a feasibility advantage.

Taken together, the experiments indicate that the main benefit of lexicographic control is not universal numerical dominance, but auditability under a transparent priority order. Scalar controls are competitive when a single compliance score is acceptable, whereas sequential $B_1$ is appropriate when the modeller wants to preserve a declared hierarchy of objectives. The joint-control, order, and lookahead checks then serve as safeguards: they test whether the sequential decomposition itself introduces quality loss beyond the chosen objective semantics.
\subsection{Reproducibility}
The public package is at \url{https://github.com/papudg/Research}, branch \texttt{reproducibility-public}, pinned to commit \path{cd25c2d90c5d48df2817c7bd2fb446355d4179b4}. It includes the metadata-only fixture, manifests, harness, dependencies, environment record, and result files.
\begin{table}[t]
\caption{Reproducibility package contents.}
\label{tab:reproducibility}
\centering
\small
\begin{tabular}{@{}>{\raggedright\arraybackslash}p{0.23\linewidth}>{\raggedright\arraybackslash}p{0.67\linewidth}@{}}
\toprule
Component & Location or execution detail \\
\midrule
Experiment harness & \texttt{experiments/}; main reproduction entry point: \texttt{python experiments/revision\_report.py} \\
Library source & \texttt{src/icaa/} \\
Dependency pins & \texttt{requirements.txt}; runs use \texttt{PYTHONHASHSEED=0} \\
Blueprint manifests & Fixed manifests under \texttt{data/manifests/} \\
Metadata fixture & \path{docs/qbank\_audit/cbse\_class\_x\_maths\_questions\_full\_updated.json}; public release excludes question text \\
Environment record & \texttt{results\_environment.json} records CPU, RAM, Windows, Python/OR-Tools versions, and fixed one-worker CP-SAT parameters \\
Run instructions & Platform-specific command list in \texttt{REPRODUCIBILITY.md} \\
\bottomrule
\end{tabular}
\end{table}
Artifact mapping: RQ1 uses \texttt{results\_revision\_suite.json}; lookahead uses \texttt{results\_lookahead\_comparison.json}; robustness uses \texttt{results\_multiseed\_revision.json}; joint/order controls use \texttt{results\_joint\_policy\_comparison.json}; tightness uses \texttt{results\_rq3\_rq5.txt}; weight sensitivity uses \texttt{results\_weight\_sensitivity.json}; and environment settings use \texttt{results\_environment.json}. Runtime entries are solver/model-construction measurements, not end-to-end preprocessing latency.
\section{Discussion and Limitations}
\label{sec:discussion}
%===============================================================================
\paragraph{Policy Design Trade-Offs.} First-feasible generation leaves structural alignment uncaptured. Multi-pass $B_1$ gives an explicit priority hierarchy---difficulty first, coverage second, Bloom fit third---and is therefore a governance choice rather than the fastest policy. In this benchmark, sequential lexicographic decomposition preserves joint-control quality, suggesting that a decomposed policy can be acceptable when the instance structure leaves enough downstream flexibility. Tighter latency budgets may instead favor $B_C$ or $B_2$; the paper's contribution is to make this trade-off visible under a common model rather than to declare a universally dominant policy. For optimization practice, the result highlights why policy semantics can matter as much as aggregate score: two policies may return similarly feasible papers while encoding different priorities and reviewability.
\paragraph{Limitations.} Structural alignment does not imply psychometric validity, and runtimes exclude embedding and conflict-graph construction. The evaluation uses one 475-item corpus and one template family, so cross-subject and cross-template generalization is untested. Difficulty/Bloom labels are source metadata, the duplicate detector was not human-validated, and fixed solver settings make sequential--joint agreement instance-specific.
These limitations also identify natural extensions: evaluating additional subjects and template families, validating duplicate edges with human review, and measuring full pipeline latency once embedding and graph construction are integrated into the benchmark.
%===============================================================================
\section{Conclusion}
\label{sec:conclusion}
%===============================================================================
We presented a policy-aware CP-SAT framework for sequential examination-paper assembly. Across five manifests, sequential lexicographic $B_1$ has median difficulty deviations of 38--40 versus 50--52 for first-feasible $B_0$. On the fixed manifest, $B_1$ and joint $B_J$ tie on all 70 comparable cases, alternative section orders retain the same aggregate medians, and bounded lookahead gives no quality gain while increasing runtime. These results support sequential lexicographic CP-SAT as an auditable policy for repeated constrained subset-selection workloads in the tested regime. More broadly, the study shows how joint controls, scalar controls, and lookahead diagnostics can be used to audit a deployed decomposition without changing the underlying CP-SAT solver. End-to-end latency, psychometric validity, and cross-corpus generalization remain future work.
The released artifacts are intended to support such extensions by allowing future work to compare new policies against the same manifests, controls, and solver settings rather than relying only on single-policy demonstrations.
\begin{thebibliography}{99}
\bibitem{vanderLinden2005}
van der Linden, W.J.: Linear Models for Optimal Test Assembly. Springer, New York (2005)

\bibitem{vanderLinden1989}
van der Linden, W.J., Boekkooi-Timminga, E.: A maximin model for IRT-based test assembly. Psychometrika 54, 185--198 (1989)

\bibitem{theunissen1985}
Theunissen, T.J.J.M.: Binary programming and test design. Psychometrika 50, 129--141 (1985)

\bibitem{lord1980}
Lord, F.M.: Applications of Item Response Theory to Practical Testing Problems. Lawrence Erlbaum, Hillsdale (1980)

\bibitem{nieuwenhuis2006}
Nieuwenhuis, R., Oliveras, A., Tinelli, C.: Solving SAT and SAT modulo theories: from an abstract Davis--Putnam--Logemann--Loveland procedure to DPLL(T). J. ACM 53, 937--977 (2006)

\bibitem{bloom1956}
Bloom, B.S. (ed.): Taxonomy of Educational Objectives, Handbook I: The Cognitive Domain. David McKay, New York (1956)

\bibitem{anderson2001}
Anderson, L.W., Krathwohl, D.R. (eds.): A Taxonomy for Learning, Teaching, and Assessing. Longman, New York (2001)

\bibitem{ehrgott2005}
Ehrgott, M.: Multicriteria Optimization, 2nd edn. Springer, Berlin (2005)

\bibitem{ignizio1976}
Ignizio, J.P.: Goal Programming and Extensions. Lexington Books, Lexington (1976)

\bibitem{perron2023}
Perron, L., Didier, F.: The OR-Tools CP-SAT Solver. Technical Report / User Manual, Google (2023)

\bibitem{hwang2006}
Hwang, G.-J.: An effective approach for test-sheet composition with large-scale item banks. Computers \& Education 46(2), 122--139 (2006)

\bibitem{moenep2020}
Ministry of Education, Government of India: National Education Policy 2020. Government of India (2020)

\bibitem{garey1979}
Garey, M.R., Johnson, D.S.: Computers and Intractability: A Guide to the Theory of NP-Completeness. W.H. Freeman, San Francisco (1979)

\bibitem{ignjatovic2021}
Ignjatovi\'c, M., Boji\'c, I., Tartalja, I.: A survey on problem formulations and (meta)heuristic-based solutions in automated assembly of parallel test forms. International Journal of Software Engineering and Knowledge Engineering \textbf{31}(08), 1171--1212 (2021). https://doi.org/10.1142/S0218194021500376

\bibitem{li2021}
Li, G., Cai, Y., Gao, X., Wang, D., Tu, D.: Automated test assembly for multistage testing with cognitive diagnosis. Frontiers in Psychology \textbf{12}, 509844 (2021). https://doi.org/10.3389/fpsyg.2021.509844

\bibitem{nguyen2021}
Nguyen, T., Bui, T., Fujita, H., Hong, T.-P., Loc, H.D., Sn\'a\v{s}el, V., Vo, B.: Multiple-objective optimization applied in extracting multiple-choice tests. Engineering Applications of Artificial Intelligence \textbf{105}, 104439 (2021). \url{https://doi.org/10.1016/j.engappai.2021.104439}

\bibitem{bennetto2021}
Bennetto, R., van Vuuren, J.H.: Multi-objective evolutionary search strategies in constraint programming. Operations Research Perspectives \textbf{8}, 100177 (2021). https://doi.org/10.1016/j.orp.2020.100177
\end{thebibliography}

\end{document}